% LaTeX template for Finance summaries
% Stu Alden
% Amortization
%%%%%%%%%%%%%%%%%%%%%%%%%%%%%%%%%%%%%%%%%%%%%%%%%%%%%%%%%%%%%%%%%%%%%%%%
\documentclass[14pt]{extarticle}

\usepackage{conceptsheet}

%---------------------------------------
\begin{document}

\finmathtitle{5. Amortization}

\textbf{Amortization} in its simplest form is the repayment of a debt with level payments ("installments"),
where each installment covers the full amount of interest that has accrued, plus a portion of the outstanding
principal.

An \textbf{amortization schedule} is a table with a line for each payment (showing the
total payment, the interest portion, and the principal portion), and the outstanding principal immediately after
that payment has been made.  At the beginning, principal is just $ A $ (present value of the annuity), and at
the end, principal goes to \$0.

To illustrate, using \$100 over 4 periods at 10\% per period:

\begin{center}
  \begin{tabularx}{\textwidth}{@{}
      l
      >{\raggedleft\arraybackslash}p{2.2cm}   % fixed width for Starting Principal
      >{\raggedleft\arraybackslash}p{2cm}
      >{\raggedleft\arraybackslash}p{2cm}
      >{\raggedleft\arraybackslash}p{2.2cm}
      >{\raggedleft\arraybackslash}p{2.5cm}
    @{}}
    \toprule
    &
    \multicolumn{1}{c}{\bfseries Starting}
    &
    \multicolumn{3}{c}{\bfseries Payment}
    &
    \multicolumn{1}{c}{\bfseries\hspace{0.5cm} Ending} \\
    \cmidrule(lr){3-5}
    \bfseries Period
    &
    \multicolumn{1}{c}{\bfseries Principal}
    &
    \multicolumn{1}{c}{\bfseries\hspace{0.2cm} Total}
    &
    \multicolumn{1}{c}{\bfseries Interest}
    &
    \multicolumn{1}{c}{\bfseries Principal}
    &
    \multicolumn{1}{c}{\bfseries\hspace{0.5cm} Principal} \\
    \midrule
    1 & \$100.00  & \$31.55 & \$10.00  & \$21.55 & \$78.45 \\
    2 &    78.45  &   31.55 &    7.85  &   24.60 &   54.75 \\
    3 &    54.75  &   31.55 &    5.47  &   26.08 &   28.67 \\
    4 &    28.67  &   31.55 &    2.87  &   28.68 &   0.00  \\
    \bottomrule
  \end{tabularx}
\end{center}

If the outstanding principal $ P $ is needed at any point $ k $ during the term of
the loan, it can be calculated using the standard annuity formulas, either \textit{prospectively} or
\textit{retrospectively}.  Here, if

\begin{description}[align=right,labelwidth=3.5cm,leftmargin=!,font=\bfseries]
  \item[n, R, i, A] As defined previously
  \item[k] Number of periods, from beginning, at which to measure outstanding principal (so $k$ is between 0 and $n$ inclusive)
  \item[P] Outstanding principal immediately after the $k$th payment
  \item[$\mathbf{\sx{\angl{k} i}}$] Accumulated value of $k$ payments thus far
  \item[$\mathbf{\ax{\angl{n-k} i}}$] Present value of the $(n-k)$ remaining payments
\end{description}

\pagebreak

Then the prospective (forward-looking) formula is
%\vspace{-.3in}

\begin{align*}
  P & = R \cdot \ax{ \angl{n-k} i} \\
\end{align*}

In words, the remaining principal is the present value of the remaining payments.

%\vspace{-.3in}

The retrospective (backward-looking) formula is

%\vspace{-.3in}

\begin{align*}
  P & = A \cdot (1+i)^k - \sx{ \angl{k} i} \\
\end{align*}

In words, the remaining principal is the accumulated value of the original loan
amount, minus the accumulated value of payments made so far.

\vspace{.2in}
\textbf{Notes:}

\begin{itemize}
  \item One reason for distinguishing between principal and interest in loan repayments is that for U.S.
    federal tax purposes, interest is (potentially) tax deductible, whereas principal is not.
  \item Knowing outstanding principal is also useful--for example, the market value of a home minus the
    outstanding principal on a mortgage represents the homeowner's \textbf{equity} in the home.
\end{itemize}

\end{document}
